\section{Model, Assumptions, and Goals}
\label{sec:model}

This section states the external conditions under which the protocol is
studied and the properties its construction must establish.  L1 and
Espresso are sources of finalized data; their internal consensus
protocols are outside the model.  The rotation, publication, and recovery
rules remain part of the protocol being specified.

\subsection{Participants, time, and failures}
\label{sec:model-participants}

% P-2.01
The system serves one L2 chain under a fixed chain identifier and rule
version.  Its registered regions are denoted by $\mathcal L_R$.  Users
submit ordinary transactions to block builders.  Stakers register
operator keys and escrow through the L1 Authority, which also records
multisignature controls.  Regional replicas run session programs; a
session $s=(\beta,r)$ is one regional turn on branch version $\beta$.
Publishers relay complete regional results to Espresso.  Batchers submit
the corresponding batch data through a trusted certifier and the L1
inbox.  External observers either follow public finality directly or
derive the L2 chain from finalized L1 inputs and DA data.  These are
logical roles: one operator may perform several of them.  Precise
identifiers and shared objects are defined in
Section~\ref{sec:shared-objects}.

Real time is denoted by $\tau$; $t$ is reserved for a leader term.  An L1
block has height $\ell$, and its transactions and events have a fixed
order.  Espresso block $J_q$ has height $q$ and timestamp $\theta_q$.
These source coordinates are distinct from L2 block height and from the
time at which a particular observer receives data.  A global clock
provides a common interpretation of local protocol deadlines without
requiring simultaneous receipt of messages or source records.

% P-2.02
The adversary may control regional participants up to the weight bound
below, send arbitrary messages, and schedule deliveries within the
applicable communication bounds.  The recovery model additionally
allows Espresso to stop finalizing new blocks.  Such an outage suspends
the healthy-period Espresso progress guarantees, while its previously
finalized history remains consistent and immutable.  Availability of
historical bytes is a separate condition.  This model does not include
conflicting finalized Espresso histories, failure of L1 finality, or a
compromised trusted certifier.

\subsection{External assumptions}
\label{sec:external-assumptions}

All service bounds in this subsection apply within the declared
admission, input-size, and load limits, with the required fee resources
and working submission routes.  Their consumers are correct
participants that remain online and keep up with the relevant history.
A request is \emph{submitted} when the complete request reaches the
service's admission boundary, rather than when a node attempts a local
send.  Bounds that include visibility cover dissemination, retrieval,
and the specified verification work as well as finalization.

% P-2.03 / AS-01
\Needspace{7\baselineskip}
\begin{assumption}[Finalized L1 history and service]
\label{ass:l1}
\textup{(S)} Authenticated finalized L1 prefixes are mutually compatible
and persistent.  They determine the same transaction order, execution
outcomes, events, and contract state at each finalized height.  A
consumer can authenticate the exact contents and completeness of a
specified range of relevant records, including an empty range.  L1
faithfully executes the deployed contract code.

\textup{(L)} There is a bound $T_{\rm L1}>0$ such that an admitted L1
operation actually submitted at time $\tau$ has a finalized execution
result by $\tau+T_{\rm L1}$, and by then the result is retrievable and
verifiable by the consumers specified above.  A finalized result may
report success or failure: state-dependent competing calls are not all
guaranteed to succeed.
\end{assumption}

The authenticated state includes Authority snapshots, bans, recovery and
closure records, the control log, and the inbox's acceptance and mode
events.  A selected collection of individually valid events does not
establish a complete control prefix.  The L1 bound begins at actual
submission and excludes fault detection, governance deliberation, and
signature collection beforehand.

% P-2.04 / AS-02
\Needspace{7\baselineskip}
\begin{assumption}[Finalized Espresso history and healthy-period service]
\label{ass:espresso}
\textup{(S)} Finalized Espresso blocks form one persistent, publicly
verifiable history.  Each block determines the complete ordered content
of this chain's namespace, including authenticated absence of content.
Its timestamps are non-decreasing.  If $J_q$ first becomes final at real
time $\tau_E(q)$, then, for a fixed $\epsilon_E\ge0$,
\begin{equation}
\label{eq:espresso-clock-lead}
 \theta_q\le\tau_E(q)+\epsilon_E.
\end{equation}

\textup{(L)} For a fixed $T_E>0$, if Espresso remains healthy throughout
$[\tau,\tau+T_E]$, an admissible complete publication actually submitted
at time $\tau$ is finalized and visible to the specified consumers by
$\tau+T_E$.  Every newly
finalized block, including an empty relevant namespace, is retrievable
and verifiable within $T_E$ of finalization under the same service
conditions.  If Espresso remains healthy, finalized block heights and
timestamps progress without bound.
\end{assumption}

Clause~(S) continues to hold in the outage model.  Clause~(L) is not used
to bound a new publication during an outage.  Finalized inclusion fixes
the published bytes and their source positions; application validity
and permission to extend the global chain are decided by this
protocol.  The timestamp condition bounds clock lead, but does not
identify a block's finalization time with its timestamp or supply a
bound on arbitrary historical retrieval.

% P-2.05 / AS-03
\Needspace{7\baselineskip}
\begin{assumption}[Authenticated synchronous communication and clock]
\label{ass:communication}
Participants share a global clock and use authenticated communication.
A complete message sent between correct, operational members of region
$i$ is delivered and processed within $\Delta_i$.  A complete message
sent between correct, operational endpoints in different regions is
delivered and processed within $\Delta_{\rm cross}$ through the mixed
peer-to-peer backbone.  These bounds apply to the declared message and
validation-dependency sizes and processing budgets; the backbone has a
working correct route under its stated fault conditions.
\end{assumption}

Messages may be duplicated or delivered out of order, and a Byzantine
sender may remain silent.  Authentication establishes the sender and
bytes, not the truth of the sender's claim.  The bounds include the
validation and state-adoption work expressly covered by a complete
message.  Any additional state reconstruction or installation has its
own preparation budget.  In particular, synchrony does not make an
unprepared recipient Ready, and a source record need not become known
to every replica at the same instant.

% P-2.06 / AS-04
\Needspace{7\baselineskip}
\begin{assumption}[Cryptography and fixed-committee faults]
\label{ass:regional-faults}
Signatures of uncompromised keys cannot be forged, and hashes bind their
canonically encoded inputs.  Session $s$ has a fixed committee $V_s$
with positive integer voting weights $w_o$, total weight
$W_s=\sum_{o\in V_s}w_o$, and fault bound $0\le F_s<W_s/2$.  The total weight
of committee keys controlled by Byzantine participants, including later
key compromises, over the period in which their
certificates may authorize new output is at most $F_s$.  Historical
certificates retain these weights; neither a later slash nor a later
configuration reweights them.
\end{assumption}

Every regional certificate requires matching signatures from distinct
members of total weight strictly greater than $F_s$.  Consequently it
contains a correct signer, and the correct members have enough total
weight to form a certificate themselves because $W_s-F_s>F_s$.
These facts do not assert honest intersection between arbitrary pairs
of certificates.  Regional safety also uses the synchronous forwarding,
timer, and recovery-lock rules.  The accepted-lifetime restriction is
essential when interpreting old signatures.  Subsequent old-key
compromise outside that lifetime is outside this analysis.  Stake
registration alone does not establish the assumed fault bound.

% P-2.07 / AS-05
\Needspace{7\baselineskip}
\begin{assumption}[DA, retention, and state preparation]
\label{ass:data-state}
A valid DA commitment binds the exact submitted payload under the fixed
DA encoding and verification rules.  Under the specified availability
and retention conditions, its complete matching bytes can be
retrieved.  Durable storage preserves committed writes.  The declared
retention arrangements provide the historical source bytes, regional
evidence, and state-reconstruction inputs that the protocol still
permits consumers to reference, or an authorized checkpoint preserving
the required validation facts.

For a selected live region, a finite preparation bound
$T_{\rm ready}$ applies only under the stated input-size, load,
availability, and native state-service conditions, and while its
assignment remains applicable.  It covers source observation, retrieval,
reconstruction, validation, and installation needed for guarded entry.
Its origin is the specified public assignment event, not a node's later
download time.
\end{assumption}

A DA certificate does not by itself guarantee permanent retrieval.
The protocol must retain the required bytes and transfer their custody
before releasing a responsible participant.  A returning node first
hydrates its required history or verifies an authorized checkpoint;
there is no uniform bound on arbitrary offline catch-up.  An observer
that cannot retrieve an already accepted payload waits for that payload
rather than replacing it with a deposit-only block.  The same exact
anchor and executable state are required after preparation, regardless
of whether they were transferred or reconstructed locally.

% P-2.08 / AS-06
\Needspace{7\baselineskip}
\begin{assumption}[Authorized and coordinated multisignature control]
\label{ass:control}
The multisignature committee acts as an honest, coordinated control
service under the authority $\mathcal K$ fixed at genesis.  It
authenticates only well-formed records for the intended chain and obeys
the protocol's publication and advance-announcement discipline.
Finalized control records are immutable.  For an ordinary Force update
actually submitted at time $\tau_{\rm submit}$, its Espresso threshold
satisfies
\begin{equation}
\label{eq:advance-announcement}
 \theta_a>\tau_{\rm submit}+T_{\rm L1}+\epsilon_E.
\end{equation}
During normal replay, at most one previously unprocessed valid ordinary
update first crosses its activation threshold in any Espresso block,
whether or not its target version will match.  This is a
separate issuance condition, not a consequence of a minimum spacing
$T_{\rm update}$ between thresholds.
\end{assumption}

By Assumptions~\ref{ass:l1} and~\ref{ass:espresso},
\eqref{eq:advance-announcement} makes an ordinary update visible before
the first block crossing its threshold can become final.  The service
has no private protected prefix.  Its authorization does not supply the
global boundary computation: applicability, the retained prefix, and
the effects of stale or overlapping controls are determined by the
protocol.  For STOP and RESUME, the committee announces an L1 activation
condition and follows the common-boundary procedure.  Honesty imposes
no response deadline, guarantees no eligible successor, and does not
exclude repeated interruptions of an otherwise serving region.

% P-2.09 / AS-07
\Needspace{7\baselineskip}
\begin{assumption}[Trusted certifier]
\label{ass:trusted-certifier}
The TEE certifier is modeled as a trusted third party $\mathcal T$ that
correctly executes its specified business rules on authenticated source
data.  The origin and integrity of every message actually emitted by
$\mathcal T$ are implicit in the model, and consumers directly trust
its typed conclusions.  A \textup{BATCH} message concerns the specified
batch inputs, DA binding, and cursor transition; a \textup{MODE}
message certifies a candidate public boundary; a \textup{CLOSURE}
message concerns an already executed closure and the stated future
authority floor.
\end{assumption}

This premise has no separate hardware, key-registration, or remote
attestation interface.  L1 contracts consume the trusted business
conclusions without verifying Espresso history on L1, and Espresso
does not verify L1 history.  A MODE certificate may precede the L1 event
that commits its boundary; a CLOSURE conclusion cannot turn a pending
announcement into an executed closure.  Correct execution by
$\mathcal T$ does not define an unspecified boundary algorithm or
establish that a proposed state root has settled.  Any required
response-time bound for the certifier is a separate progress condition.

% P-2.10 / AS-08
\Needspace{7\baselineskip}
\begin{assumption}[Native execution and state-settlement soundness]
\label{ass:native-execution}
The pinned native rule version deterministically constructs and executes
an L2 block from its exact parent state, ordered ordinary transactions,
and the required finalized L1 context, including the origin header,
deposits, system inputs, and configuration.  Native validation and
state-adoption checks are sound for those inputs.  Under its declared
verification and challenge assumptions, the existing state-settlement
mechanism finalizes only state claims consistent with the canonical
derived inputs and their prescribed execution.
\end{assumption}

The manuscript uses this boundary for MegaETH's execution and
Kailua/fault-proof path; it does not specify a replacement execution or
dispute system.  A transaction that executes and reverts may be part of
a valid block.  Input correspondence certified by $\mathcal T$ and
correctness of the resulting state claim remain distinct statements.
In particular, sound execution of a derived deposit-only block would
not justify substituting it for a different block already confirmed in
$G$.  Establishing correspondence with $G$ is a protocol obligation.

\begin{table}[htbp]
\centering\small
\caption{External assumptions and the conditions used by the protocol.
Healthy-period service clauses are not premises about an Espresso
outage.}
\label{tab:assumption-coverage}
\begin{tabularx}{\linewidth}{@{}P{.16\linewidth}P{.27\linewidth}P{.26\linewidth}Y@{}}
\toprule
Assumption & Safety or interpretation premise & Availability or timing
condition & Main consumers\\
\midrule
L1 (\ref{ass:l1}) & Consistent finalized state, events, and complete
ranges & Admitted-operation result and visibility within $T_{\rm L1}$
& Authority, replay, inbox, derivation\\
Espresso (\ref{ass:espresso}) & Persistent ordered history, complete
namespace, timestamp bound & $T_E$ and ongoing header/time progress
while healthy & Publication and NORMAL replay; retained history in recovery\\
Network (\ref{ass:communication}) & Authenticated messages and a common
clock & Regional $\Delta_i$, cross-region $\Delta_{\rm cross}$ within
the complete-message budget & Programs, Managers, source relay\\
Faults (\ref{ass:regional-faults}) & Binding hashes, unforgeability, fixed
weight bound $F_s<W_s/2$ & Operational participation is required in
addition to honest weight & Certificates and accountability\\
Data (\ref{ass:data-state}) & Exact DA binding and durable writes &
Retention, retrieval, and scoped $T_{\rm ready}$ & Publishing, state
preparation, batching, derivation\\
Control (\ref{ass:control}) & Authorized immutable records and announcement
discipline & No response or uninterrupted-service bound follows from
honesty & Force and mode control\\
Certifier (\ref{ass:trusted-certifier}) & Trusted, typed business
conclusions & Any response bound must be stated separately & Inbox,
mode submission, clearance\\
Native rules (\ref{ass:native-execution}) & Deterministic execution and
sound state settlement & Input availability and settlement progress
have their own conditions & Builders, replicas, observers, state claims\\
\bottomrule
\end{tabularx}
\end{table}

The Authority transitions, session rules, publication duties, and common
mode-boundary procedure are constructions to be specified under these
assumptions.  In particular, neither prefix stability of the boundary
policy nor preservation of globally confirmed inputs by the settlement
path is an assumed property of L1 or Espresso.  A concrete protection
profile must establish that connection through its delivery conditions
or its explicit inbox and derivation rules.

\subsection{Protocol goals}
\label{sec:protocol-goals}

% P-2.11
\paragraph{Safety and accountability.}
Let $G_m^u$ be the global block sequence confirmed by correct observer
$u$ after its $m$th durable replay step, and let $\preceq$ denote the
prefix relation on exact block sequences.  The public rule should
establish
\begin{align}
\label{eq:goal-agreement}
 &G_m^u\preceq G_n^v\quad\text{or}\quad G_n^v\preceq G_m^u
 &&\text{for all correct }u,v\text{ and steps }m,n,\\
\label{eq:goal-persistence}
 &G_m^u\preceq G_{m+1}^u.
\end{align}
Every non-genesis block also retains the producing context or L1-derived
provenance recorded when it was appended.  Let $\widehat G$ be a
correctly derived native chain from the same authorized starting point.
The stronger settlement goal is that, whenever $B\in G_m^u$ and
$\widehat G$ has reached height $h(B)$,
\begin{equation}
\label{eq:goal-native-correspondence}
 \widehat G[h(B)]=\operatorname{native}(B).
\end{equation}
An append-only replay implementation alone does not establish
\eqref{eq:goal-native-correspondence}.  The protocol must additionally
keep session core states isolated, prevent repeated initial entry or
revival after termination, apply penalties only to the specified valid
offenses without duplicate deductions, and release stake only after its
remaining responsibilities have been closed and the applicable evidence
window has elapsed.

% P-2.12
\paragraph{Conditional progress.}
Progress is studied separately in three phases.  In NORMAL, suitable
regional leadership and participation, available proposal inputs,
preparation, complete-result publication, and healthy external services
should allow sessions to complete and $G$ to grow.  A finite session
aims to obtain its terminal receipt; unbounded DATA growth is a property
of successive sessions, not of one fixed quota.  After STOP, progress
depends on a committed common boundary, delivery of the protected
prefix, and available L1 and native derivation inputs.  After RESUME,
ordinary production further requires a common derived anchor, usable
Espresso service, an eligible live region, and guarded entry with the
correct state.  Any finite recovery bound must include the control,
certifier, data, and preparation services it uses.

% P-2.13
\paragraph{Scope.}
The protocol does not promise fair service among regions, inclusion of
every ordinary transaction, progress under arbitrary repeated controls,
or lossless recovery from arbitrary data loss.  Transaction-level
liveness would require an additional builder and input-fairness
condition.  L1-derived deposits and forced-inclusion transactions retain
their native treatment; they do not add a regional ordering-fairness
mechanism.  The safety statements distinguish the accepted prefix $G$
from the production branch ahead of it: a closed, unaccepted suffix has
no post-closure regional agreement or completion guarantee.  These are
goals for the specified composition and its stated conditions, rather
than a claim of universal composability or a proof of the underlying
chains' implementations.
